\section{Extraction with an inverse-genus error}
\label{sec:sharp-extraction}

The local limit theorem \eqref{eq:TL47} has no rate. Its quantitative form, Proposition~\ref{prop:quantitative-local} in Appendix~\ref{app:extraction}, gives only a power rate in $\min(h,\mu)$, and its Gaussian approximation degenerates when the coefficient index stays bounded. To obtain an error of order $g^{-1}$ uniformly in the index, we retain an exact Poisson factor rather than its Gaussian approximation. Theorem~\ref{thm:sharp-extraction} supplies the estimate used for the leading absolute-volume formula in Section~\ref{sec:absolute}, and Theorem~\ref{thm:canonical-expansion} extends it to all orders. Both theorems are proved in Appendix~\ref{app:extraction}.

\emph{Why a Poisson factor.} Near zero cusp weight the leading term of $\Phi_g^{(p)}$ behaves like $\mathrm e^{hS(x)}$ with $S(x)=2x+O(x^2)$; see~\eqref{eq:CE1} below. At a positive weight $x$, the normalized terms $\mathrm e^{-2hx}(2hx)^k/k!$ are exactly Poisson probabilities with mean $2hx$. Hence, when the coefficient index $k$ stays bounded, the cusp count under the saddle weight is approximately Poisson, and a Gaussian approximation of its point mass makes a relative error of order $1/k$, which does not tend to zero. If the leading term is written instead with the exact Poisson mass $\mathfrak p_k=\mathrm e^{-k}k^k/k!$ at the mean, the ratio of the true coefficient to this model is $1+O(1/h)$ for every $k$. For large $k$ one has $\sqrt{2\pi k}\,\mathfrak p_k\to1$, and the Gaussian formula is recovered.

Fix $p\ge1$, put $F_h=\Phi_g^{(p)}$, $h=2g-2+p$, and write
\begin{equation}
 \mathcal J_{g,p}(x)=\mathcal J_{g,p}(0)\mathrm{e}^{hS(x)}B_p(x),\qquad
 S=\log\frac{\mathcal Z(0)}{\mathcal Z(Y(x))},\qquad
 B_p=\frac{r^{p/2}\sqrt{A\Theta}}{r(0)^{p/2}\sqrt{A(0)\Theta(0)}}.
 \label{eq:CE1}
\end{equation}
We view $\mathcal J$ as a function of $x=X(Y)$. Then
$\mathscr DS=\Psi$ and $\mathscr D\log B_p=\beta_p\Psi$.
For each integer $k\ge1$, choose $x$ by $h\Psi(x)=k$, and put
$v=\mathscr D\Psi$, $\sigma^2=hv$, $M=\min(h,k)$.
Let $\mathfrak p_k=\mathrm{e}^{-k}k^k/k!$, with $\mathfrak p_0=1$.

\begin{theorem}[Uniform Poisson-corrected extraction]\label{thm:sharp-extraction}
For fixed $p\ge1$, uniformly over integers $k\ge0$,
\begin{equation}
 [z^k]F_h(z)=\frac{\mathcal J_{g,p}(x)}\pi x^{-k}\mathfrak p_k
       \sqrt{\frac{k}{h\mathscr D\Psi(x)}}\,(1+O_p(h^{-1})).
 \label{eq:CE2}
\end{equation}
At $k=0$ take $x=0$ and interpret both $k$-dependent factors by
continuity as one.
\end{theorem}

The proof, given in Subsections~\ref{app:local-complex}--\ref{app:small-indices}, has three parts: transfer a sufficiently long marked expansion, namely Theorem~\ref{thm:full-expansion} at a fixed finite order, to small complex neighborhoods of the real saddle point; identify the Poisson correction; and estimate the small indices directly. The argument uses no estimate for a growing number of markings and never differentiates a remainder; away from zero, the model is continued analytically only to small neighborhoods of the real axis.

\subsection{The full expansion at a prescribed cusp count}\label{sec:canonical-expansion}

Theorem~\ref{thm:sharp-extraction} recovers the leading term with an inverse-genus error. We now extract the whole marked expansion, keeping the remainder uniform in both the cusp count and the truncation order. The result is Theorem~\ref{thm:canonical-expansion}, which provides the remaining input for the full absolute-volume expansion in Section~\ref{sec:absolute}.

\emph{The idea.} Away from zero cusp weight, including near the critical point, the coefficient integral has a nondegenerate saddle point, and the Gaussian method of steepest descent gives an expansion in $1/h$. Near zero this fails for bounded $k$: the saddle point of $\mathrm e^{hS(z)}z^{-k}$ approaches the singularity of $z^{-k}$ at the origin. Following Temme, we change variables so that the phase $S(z)-(k/h)\log z$ becomes exactly $u-(k/h)\log u$. The coefficient is then a Poisson coefficient $[u^k]\mathrm e^{hu}f(u)$ with an analytic amplitude $f$. Expanding $f$ about $u=k/h$ produces the signed central polynomials $T_m(k)$ of the coefficient operator in~\eqref{eq:CG8}, divided by powers of $h$. These are not the central moments of a probability distribution: for example, $T_2(k)=-k$. The resulting expansion in $1/h$ is uniform in $k$, including $k=0$. The two expansions are identified on an overlap of their ranges.

The proof, given in Subsections~\ref{app:phase-change}--\ref{app:canonical-proof}, uses a Poisson integral near zero cusp weight and a Gaussian integral away from zero, including near criticality. A change of phase identifies the two coefficient formulas on their overlap. Near zero we follow Temme's method for uniform Laplace and loop integrals~\cite[Sections~2.3 and~3.7]{Temme}.

Retain $h=2g-2+p$, $k=h\Psi(x)$, $v=\mathscr D\Psi(x)$ and the
functions $S,B_p$ in \eqref{eq:CE1}. Regard $\mathcal L_{p,j}$
as functions of $x$ via $Y(x)$. Write
\begin{equation}
 \mathcal Q_{g,p,k}=
 \frac{\mathcal J_{g,p}(x)}\pi x^{-k}\mathfrak p_k
                   \sqrt{\frac{k}{hv}}.
 \label{eq:CG1}
\end{equation}
The conventions at $k=0$ are those of Theorem~\ref{thm:sharp-extraction}.


Put $\mathcal L_{p,0}=1$. Let $\mathfrak G_w$ be the linear functional
on polynomials which maps $w^{2a}$ to $(-1)^a(2a-1)!!$ and odd
powers to zero. Define the formal even series
\begin{equation}
\begin{split}
 \mathcal N_p(\tau,x)=\mathfrak G_w\Bigg[
 &\exp\!\left\{
 \sum_{m\ge3}\frac{\tau^{m-2}\mathscr D^mS}{m!v^{m/2}}w^m
 +\sum_{m\ge1}\frac{\tau^m\mathscr D^m\log B_p}{m!v^{m/2}}w^m
 \right\}\\[-2pt]
 &\hspace{20pt}\cdot\sum_{s\ge0}\tau^{2s}
       \mathcal L_{p,s}(x\mathrm{e}^{\tau w/\sqrt v})\Bigg].
\end{split}
 \label{eq:CG2}
\end{equation}
Formal composition in the last line means Taylor expansion at $x$.
At every order only finitely many derivatives are used. If
$\sqrt{2\pi k}\,\mathfrak p_k\sim\sum_{j\ge0}s_jk^{-j}$, $s_0=1$, set
\begin{equation}
 \mathcal A_{p,j}(Y)=[\tau^{2j}]
       \frac{\mathcal N_p(\tau,X(Y))}
            {\displaystyle\sum_{i\ge0}s_i\Psi(X(Y))^{-i}\tau^{2i}}.
 \label{eq:CG3}
\end{equation}
This initially defines the functions for $0<Y<y_0$.

\begin{theorem}[Uniform expansion at a prescribed cusp count]\label{thm:canonical-expansion}
For every fixed $p\ge1$ there are constants $C_p,\mathfrak B_p,c_p,h_p>0$
such that, for $h=2g-2+p\ge h_p$, every integer $k\ge0$ and
$0\le K\le c_p\sqrt h$,
\begin{equation}
 \left|\frac{[z^k]\Phi_g^{(p)}}{\mathcal Q_{g,p,k}}
       -\sum_{j=0}^K\frac{\mathcal A_{p,j}(Y)}{h^j}\right|
 \le\frac{C_p\mathfrak B_p^{K+1}((K+1)!)^2}{h^{K+1}},
 \qquad h\Psi(X(Y))=k,
 \label{eq:CG4}
\end{equation}
where $\mathcal A_{p,0}=1$.
The coefficients are analytic on one complex neighborhood $\Omega_p$
of $[0,y_0]$, with
$\|\mathcal A_{p,j}\|_{\Omega_p}\le C_p\mathfrak B_p^j(j!)^2$, and
$\mathcal A_{p,j}(0)=\mathcal L_{p,j}(0)$.
For a sufficiently small fixed $c'_p>0$, truncation at
$K=\lfloor c'_p\sqrt h\rfloor$ has relative error
$O_p(\mathrm{e}^{-a_p\sqrt h})$, uniformly in $k$, for some $a_p>0$.
The marking count is fixed. The infinite series need not converge.
\end{theorem}

The first coefficient $\mathcal A_{p,1}$ is given explicitly in \eqref{eq:CG12}, and its value at the critical endpoint in \eqref{eq:CG13}. Retaining it improves the relative remainder to $O_p(h^{-2})$ over every coefficient index.

